\documentclass{article}
\usepackage[utf8]{inputenc}
\usepackage{amsmath}
\usepackage{helvet}
\renewcommand{\familydefault}{\sfdefault}
\begin{document}

\section*{Teoria: Ecuaciones Lineales}

\textbf{Forma general}

{\Large \[ ax + b = 0 \quad \Rightarrow \quad x = \frac{-b}{a} \] }

\begin{itemize}
\item {\Large \(a\)}: coeficiente del termino con \(x\)
\item {\Large \(x\)}: la incognita a despejar
\item {\Large \(b\)}: termino independiente
\end{itemize}

Una ecuación lineal es aquella donde la incógnita aparece elevada solo a la
primera potencia, sin multiplicarse por sí misma ni aparecer dentro de raíces.
Su forma general es \(ax + b = 0\), con \(a \neq 0\).

Resolverla consiste en despejar x usando operaciones inversas: se resta b
de ambos lados, y luego se divide por a, manteniendo la igualdad en todo momento.

\subsection*{Ejemplo 1: Caso directo}
Resolver \(3x + 6 = 0\).

\[ 3x = -6 \quad \Rightarrow \quad x = -2 \]

La solución es \(x = -2\).

\subsection*{Ejemplo 2: Incógnita en ambos lados}
Resolver \(5x - 3 = 2x + 9\).

\[ 5x - 2x = 9 + 3 \quad \Rightarrow \quad 3x = 12 \quad \Rightarrow \quad x = 4 \]

La solución es \(x = 4\).

\subsection*{Ejemplo 3: Con fracciones}
Resolver \(\frac{x}{2} + 3 = 7\).

\[ \frac{x}{2} = 4 \quad \Rightarrow \quad x = 8 \]

La solución es \(x = 8\).

\subsection*{Ejemplo 4: Sin solución (caso especial)}
Resolver \(2x + 5 = 2x + 8\).

\[ 2x - 2x = 8 - 5 \quad \Rightarrow \quad 0 = 3 \]

Esta igualdad es falsa, por lo tanto la ecuación no tiene solución.

\end{document}
